\documentclass[11pt,a4paper]{article}
\usepackage[utf8]{inputenc}
\usepackage[T1]{fontenc}
\usepackage[french]{babel}
\usepackage{amsmath,amssymb,amsfonts}
\usepackage{geometry}
\usepackage{enumitem}
\usepackage{booktabs}

\geometry{top=2cm, bottom=2cm, left=2cm, right=2cm}

\title{\textbf{Progression de Mathématiques \\ Classes de 5\textsuperscript{ème} (2026)}}
\author{}
\date{}

\begin{document}

\maketitle

\section*{Trimestre 1}

% --- SÉQUENCE 1 ---
\subsection*{Séquence 1 : Décimaux et Opérations}

\paragraph{Objectifs d'apprentissage}
\begin{itemize}[label=\bullet]
    \item Rappels sur les nombres décimaux.
    \item Nommer un calcul, distinguer sommes et produits, termes et facteurs.
    \item Enchaîner des opérations.
    \item Connaître et utiliser les priorités opératoires.
    \item Additionner, soustraire, multiplier et diviser pour résoudre des problèmes et contrôler la vraisemblance de son résultat.
    \item Traduire un problème, une succession donnée d'opérations, un programme de calcul, en une seule expression, en faisant appel ou non à des parenthèses.
    \item Connaître le sens et les situations d'emploi de ces opérations.
\end{itemize}

\paragraph{Automatismes}
\begin{itemize}[label=\bullet]
    \item Reconnaître et construire le symétrique d'une figure par symétrie axiale, dont l'axe est vertical, horizontal, ou en diagonale sur quadrillage.
    \item Construire le symétrique, par rapport à un axe, d'un point, d'une figure, sur feuille blanche.
\end{itemize}


% --- SÉQUENCE 2 ---
\subsection*{Séquence 2 : Symétrie centrale}

\paragraph{Objectifs d'apprentissage}
\begin{itemize}[label=\bullet]
    \item Définir le demi-tour, ou symétrie centrale.
    \item Connaître les propriétés du demi-tour.
\end{itemize}

\paragraph{Prolongements possibles : mises en perspective historiques et culturelles}
\begin{itemize}[label=\bullet]
    \item Rosaces et pavages dans l'art du Moyen Âge.
    \item Lien avec les formes vues dans la nature : flocon, alvéoles d'abeilles, fleurs.
\end{itemize}

\paragraph{Automatismes}
\begin{itemize}[label=\bullet]
    \item Placer sur une demi-droite graduée un point dont l'abscisse est un nombre décimal.
    \item Repérer un nombre décimal sur une demi-droite graduée.
\end{itemize}


% --- SÉQUENCE 3 ---
\subsection*{Séquence 3 : Nombres relatifs (1)}

\paragraph{Objectifs d'apprentissage}
\begin{itemize}[label=\bullet]
    \item Définir les nombres relatifs.
    \item Définir l'opposé et la valeur absolue d'un nombre.
    \item Définir la notion de nombre positif, strictement positif, négatif, strictement négatif.
    \item Utiliser les nombres relatifs pour représenter des grandeurs observables qui peuvent prendre des valeurs inférieures à zéro (température, temps, altitude, etc.), en particulier dans le cadre de la résolution de problèmes.
    \item Lire l'abscisse d'un nombre relatif sur une droite graduée et placer un nombre relatif d'abscisse donnée.
    \item Comparer et ranger dans l'ordre croissant et décroissant des nombres décimaux relatifs.
    \item Dans le plan muni d'un repère orthogonal :
    \begin{itemize}[label=-]
        \item lire les coordonnées d'un point donné ;
        \item placer un point de coordonnées données.
    \end{itemize}
\end{itemize}

\paragraph{Prolongements possibles : mises en perspective historiques et culturelles}
\begin{itemize}[label=\bullet]
    \item Premières apparitions des nombres négatifs avec Brahmagupta (600), dont l'existence était toujours niée au XVII\textsuperscript{e} siècle.
\end{itemize}

\paragraph{Automatismes}
\begin{itemize}[label=\bullet]
    \item Reconnaître et citer sur une configuration géométrique des angles : angle plein, plat, nul, droit ; angles opposés par le sommet, adjacents, supplémentaires, aigus, obtus.
    \item Savoir qu'un angle droit mesure 90°, qu'un angle plat mesure 180°.
    \item Reconnaître une bissectrice.
    \item Connaître les mesures des angles de l'équerre dont dispose l'élève (30° 60° 90° ou 45° 45° 90°).
    \item Reconnaître un triangle isocèle, équilatéral ou rectangle à partir d'un schéma codé.
\end{itemize}


% --- SÉQUENCE 4 ---
\subsection*{Séquence 4 : Parallélisme et angles}

\paragraph{Objectifs d'apprentissage}
\begin{itemize}[label=\bullet]
    \item Caractériser le parallélisme par les angles : angles alternes internes, angles correspondants.
\end{itemize}

\paragraph{Prolongements possibles : mises en perspective historiques et culturelles}
\begin{itemize}[label=\bullet]
    \item Éléments d'Euclide (Livre I, proposition 27 ou 28).
\end{itemize}

\paragraph{Automatismes}
\begin{itemize}[label=\bullet]
    \item Savoir que pour compléter une addition à trou, on utilise une soustraction.\\
    \textit{Exemple :} $2 + \dots = 7$ se complète en calculant $7 - 2$.
    \item Additionner et soustraire des décimaux.\\
    \textit{Exemples :} $2,7 + 1,4$ ; $3,4 - 0,8$.
    \item Additionner, soustraire, multiplier des nombres décimaux à une ou deux décimales.
\end{itemize}


% --- SÉQUENCE 5 ---
\subsection*{Séquence 5 : Calcul littéral et algébrique (1)}

\paragraph{Objectifs d'apprentissage}
\begin{itemize}[label=\bullet]
    \item Produire des formules (double, triple, carré, successeur, prédécesseur, aire, périmètre, etc.).
    \item Réduire une expression littérale de la forme $a + b$, où $a$ et $b$ sont des nombres décimaux.
    \item Calculer la valeur d'une expression littérale par substitution.
    \item Tester si une égalité entre expressions algébriques comportant une variable est vraie ou fausse.
    \item Utiliser un contre-exemple pour démontrer qu'une assertion est fausse.
\end{itemize}

\paragraph{Automatismes}
\begin{itemize}[label=\bullet]
    \item Reconnaître si une situation donnée entre dans le cadre de la proportionnalité ou non.
    \item Dans des situations simples, mobiliser une procédure adaptée (propriété de linéarité pour la multiplication ou l'addition, retour à l'unité) pour résoudre un problème lié à la proportionnalité.\\
    \textit{Exemples :}
    \begin{itemize}[label=-]
        \item à partir d'une recette pour 4 personnes, on sait donner (ou verbaliser la procédure) les quantités lorsque l'on passe à 2, 8 ou 6 personnes.
        \item si l'on connaît le prix d'un kilogramme de tomates, on sait comment calculer le prix de 3 kg ou de 4,3 kg de tomates.
        \item lors de l'élection des délégués de la classe, 4 élèves se présentent. Chaque élève a voté pour un seul candidat. Voici les résultats :
    \end{itemize}
\end{itemize}

\begin{center}
\begin{tabular}{lcccc}
    \toprule
    \textbf{Candidat} & Alice & Philippe & Chloé & Dany \\
    \midrule
    \textbf{Nombre de voix} & 10 & 7 & 5 & 3 \\
    \bottomrule
\end{tabular}
\end{center}

\begin{itemize}[label=\bullet]
    \item[] Calculer le pourcentage de voix de chaque candidat.
\end{itemize}


% --- SÉQUENCE 6 ---
\subsection*{Séquence 6 : Triangles (1)}

\paragraph{Objectifs d'apprentissage}
\begin{itemize}[label=\bullet]
    \item Construire des triangles à partir de données partielles.
    \item Connaître la somme des angles d'un triangle et savoir la démontrer.
    \item Utiliser ces propriétés dans le cas de triangles particuliers.
    \item Définir et tracer les hauteurs dans un triangle.
    \item Savoir que les hauteurs d'un triangle sont concourantes.
    \item Résoudre des problèmes faisant appel à des conversions d'unités de longueur et d'unités d'aires.
    \item Calculer l'aire d'un triangle.
\end{itemize}

\paragraph{Automatismes}
\begin{itemize}[label=\bullet]
    \item Mobiliser les critères de divisibilité par 2, 5 et 10 vus en CM1 et CM2.
    \item Déterminer le quotient et le reste dans une division euclidienne.\\
    \textit{Exemple :} $17 = 3 \times 5 + 2$.
\end{itemize}

\newpage

\section*{Trimestre 2}

% --- SÉQUENCE 7 ---
\subsection*{Séquence 7 : Nombres rationnels (1)}

\paragraph{Objectifs d'apprentissage}
\begin{itemize}[label=\bullet]
    \item Formes d'un rationnel.
    \item Égalité de quotient : simplification des fractions rationnelles.
    \item Application à la division avec un diviseur décimal.
    \item Repérer et placer une fraction sur un axe gradué.
    \item Comparaison de fractions.
    \item Encadrement de fractions entre deux entiers consécutifs.
    \item Diviser par un nombre décimal.
\end{itemize}

\paragraph{Automatismes}
\begin{itemize}[label=\bullet]
    \item Comparer deux fractions.\\
    \textit{Exemples :} $\dfrac{2}{7}$ et $\dfrac{5}{7}$ ; $\dfrac{8}{12}$ et $\dfrac{8}{21}$ ; $\dfrac{3}{4}$ et $\dfrac{7}{18}$ ; $\dfrac{8}{3}$ et $\dfrac{6}{7}$.
    \item Écrire une fraction sous la forme d'une somme d'un nombre entier et d'une fraction inférieure à 1.\\
    \textit{Exemple :} $\dfrac{17}{5} = 3 + \dfrac{2}{5}$.
\end{itemize}


% --- SÉQUENCE 8 ---
\subsection*{Séquence 8 : Triangles (2)}

\paragraph{Objectifs d'apprentissage}
\begin{itemize}[label=\bullet]
    \item Connaître les propriétés des médiatrices et du cercle circonscrit dans le cas de triangles particuliers.
    \item Définir et tracer les médianes dans un triangle.
    \item Démontrer qu'une médiane partage un triangle en deux triangles d'aires égales.
\end{itemize}

\paragraph{Prolongements possibles : mises en perspective historiques et culturelles}
\begin{itemize}[label=\bullet]
    \item Droite d'Euler, cercle des neufs points.
    \item En suivant un programme de construction donné, l'élève construit, à la règle et au compas, un pentagone régulier.
    \item Construction d'ogives en lien avec l'art gothique.
    \item Construction de rosaces trilobées, du triangle de Reuleaux.
    \item Différents instruments de mesure.
\end{itemize}

\paragraph{Automatismes}
\begin{itemize}[label=\bullet]
    \item Connaître la notion de médiatrice, de cercle circonscrit.
\end{itemize}


% --- SÉQUENCE 9 ---
\subsection*{Séquence 9 : Nombres relatifs (2)}

\paragraph{Objectifs d'apprentissage}
\begin{itemize}[label=\bullet]
